% Moved from 02-line.tex (the length take, question (a)): the line statements that Section 3
% proves again at d = 1, with their proofs of record.
\section{The strata on the line}
\label{supp:strata}

This section states and proves on the line the statements that Section~\ref{ch:pyramid} of the
paper proves in $\RR^d$: the closed subgroups of dilations, discrete covariance in \Levy{}
coordinates, the three strata, the two families that separate them, and the absence of lattice
kernels. They are written here in the line's own terms, the operator axioms (A1)--(A8) and the
folded data $(a(t), \varpi_t)$, because Section~\ref{ch:line-appendix} and some proofs of the paper
read them in that notation. Section~\ref{subsec:line-strata} of the paper says which statement of
Section~\ref{ch:pyramid} each of them is at $d = 1$.

Write $\mathrm{ND}_s$ for the set of functions
\[
  \psi(\omega) = a\,\omega^2 + \int_{(0,\infty)}\bigl(1 - \cos\omega v\bigr)\,\varpi(dv), \qquad
  a \ge 0, \quad \int(1 \wedge v^2)\,\varpi(dv) < \infty ,
\]
and call $(a, \varpi)$ the \emph{data} of $\psi$: a Gaussian coefficient and a folded \Levy{}
measure. The first lemma is what the supplement uses of the \Levy{}--Khintchine formula.

% shared with blueprint lem:line-exponent-data
\begin{lemma}[the data of an exponent]
\label{lem:line-exponent-data}

\begin{enumerate}
\item Every $\psi \in \mathrm{ND}_s$ is $-\log\widehat\mu$ for exactly one symmetric infinitely
divisible law $\mu$ on $\RR$, and every such $-\log\widehat\mu$ lies in $\mathrm{ND}_s$. The data of
$\psi$ are unique.
\item If $\psi$ has the data $(a, \varpi)$ and $\lambda > 0$, then $\psi(\lambda\,\cdot)$ has the
data $(\lambda^2a, T_\lambda\varpi)$.
\item If $\psi, \psi' \in \mathrm{ND}_s$ have the data $(a, \varpi)$ and $(a', \varpi')$, then
$\psi' - \psi \in \mathrm{ND}_s$ iff $a' \ge a$ and $\varpi' \ge \varpi$, and then $\psi' - \psi$
has the data $(a' - a, \varpi' - \varpi)$.
\end{enumerate}
\end{lemma}

\begin{proof}
(1) Let $\nu$ be a symmetric \Levy{} measure on $\RR$. The imaginary part of the integrand,
$\sin\omega x - \omega x\mathbf 1_{|x| \le 1}$, is odd and $\nu$-integrable, so its integral
vanishes, and
\[
  \int\bigl(e^{i\omega x} - 1 - i\omega x\mathbf 1_{|x| \le 1}\bigr)\,\nu(dx)
  = -\int(1 - \cos\omega x)\,\nu(dx).
\]
So the law with the triplet $(2a, \nu, 0)$, $\nu$ the symmetrization $\tfrac12(\varpi + \varpi(-\,\cdot))$
of $\varpi$, has the exponent $\psi$. It exists by [\sscite{sato1999levy}, Thm.~8.1(iii),
p.~38]. It is symmetric because its transform $e^{-\psi}$ is real and even. A symmetric law
has a symmetric triplet, by the uniqueness of the triplet [\sscite{sato1999levy}, Thm.~8.1(ii),
p.~38] applied to the law and its
reflection, so its exponent is of this form. The law with exponent $\psi$ is unique because the
transform determines the law. The data are unique because the triplet is.

(2) $\int(1 - \cos\lambda\omega v)\,\varpi(dv) = \int(1 - \cos\omega u)\,(T_\lambda\varpi)(du)$ by
the substitution $u = \lambda v$, and $a(\lambda\omega)^2 = \lambda^2a\,\omega^2$. Then (1) gives
the uniqueness.

(3) Let $a' \ge a$ and $\varpi' \ge \varpi$. The difference $\varpi' - \varpi$ is a measure: $\varpi$
is finite on each of the disjoint sets $\{v \le 0\} \cup \{v > 1\}$ and
$\{1/(n+1) < v \le 1/n\}$, $n \ge 1$, which cover $\RR$, and $\varpi' - \varpi$ is the sum of the
differences of the restrictions to these sets. It is at most $\varpi'$, so it is a \Levy{} measure
on $(0,\infty)$, and $\psi' - \psi$ has the data $(a' - a, \varpi' - \varpi)$ by linearity of the
formula. Conversely, if $\psi' - \psi$ has data $(a'', \varpi'')$, then $\psi' = \psi + (\psi' - \psi)$
has the data $(a + a'', \varpi + \varpi'')$, which are $(a', \varpi')$ by (1).

Every step is machine-checked, and the cited fact is \ledger{A39}, read on $\RR$ inside $\RR^1$
through $x \mapsto xe_1$.
\end{proof}

% shared with blueprint lem:closed-subgroups
\begin{lemma}[closed subgroups]
\label{lem:closed-subgroups}

The closed subgroups of $\bigl((0,\infty), \times\bigr)$ are $\{1\}$, the cyclic groups
$q^{\ZZ}$ with $q > 1$, and $(0,\infty)$. A family satisfying (A1)--(A7), (ND), and (A8) for
every dilation in a set $\Lambda \subset (0,\infty)$, satisfies (A8) for every dilation in the
closed subgroup that $\Lambda$ generates.
\end{lemma}

\begin{proof}
The logarithm is an isomorphism of $\bigl((0,\infty),\times\bigr)$ onto $(\RR,+)$ and a
homeomorphism. Let $H \ne \{0\}$ be a closed subgroup of $\RR$. Since $H = -H$, it has positive
elements; let $\ell$ be the infimum of its positive elements. Suppose $\ell > 0$. Since $H$ is
closed, $\ell \in H$, so $\ell$ is the least positive element of $H$. For $x \in H$ let $k$ be
the integer with
$k\ell \le x < (k+1)\ell$. Then $x - k\ell \in H$ lies in $[0, \ell)$, so it is $0$, and
$H = \ell\ZZ$. Suppose $\ell = 0$. Then for every $\delta > 0$ the group $H$ has an element in
$(0,\delta)$, whose multiples meet every interval of length $\delta$. So $H$ is dense and, being
closed, equal to $\RR$. With $q := e^\ell$ this gives the three groups. Conversely $\{1\}$ and
$(0,\infty)$ are closed subgroups, and so is $q^{\ZZ}$ for $q > 1$: its logarithms are the
multiples of $\log q$, a discrete and hence closed subgroup of $\RR$.

For the second clause, the kernels $\mu_{s,t}$ of the family are a symmetric family in dimension
$1$, and (A8) at a dilation $\lambda$ is $\lambda$-covariance of the kernels
(Lemma~\ref{lem:line-axioms-kernels}). By Lemma~\ref{lem:family-dilations}(2) the set $G$ of
$\lambda$ at which the kernels are $\lambda$-covariant is a closed subgroup of $(0,\infty)$. It
contains $\Lambda$, so it contains the closed subgroup that $\Lambda$ generates, and at each of
its elements Lemma~\ref{lem:line-axioms-kernels}(3) turns $\lambda$-covariance back into (A8).

Every step is machine-checked, and the cited facts enter only through
Lemma~\ref{lem:family-dilations}(2).
\end{proof}

So the dilations under which a family is covariant form $\{1\}$, some $q^{\ZZ}$ or $(0,\infty)$.
The strata of Theorem~\ref{thm:lattice-line} below are nested: the first is every family with
(A1)--(A7) and (ND), the second those covariant under the powers of one ratio $q$, the third the
fully covariant ones. In the second, Lemma~\ref{lem:family-dilations}(3) at $d = 1$ relabels the scales so that $S_qt = qt$, the
canonical discrete gauge. That gauge is unique only up to the group $\mathcal{G}_q$ of
Lemma~\ref{lem:gauge-group}.

% shared with blueprint lem:discrete-covariance-line
\begin{lemma}[discrete covariance in \Levy{} coordinates]
\label{lem:discrete-covariance-line}

Let $(\Phi_{s,t})$ satisfy (A1)--(A7) and (ND), with kernels $\mu_{s,t}$
(Lemma~\ref{lem:line-axioms-kernels}). Let $\Psi_t := -\log\widehat\mu_{0,t}$, and let
$(a(t), \varpi_t)$ be the folded data of $\mu_{0,t}$: with the \Levy{} data $(A_t, \nu_t)$ of
Lemma~\ref{lem:family-structure} at $d = 1$, $a(t) := A_t$, a number since $d = 1$, and
$\varpi_t$ is the image of $\nu_t$ under $x \mapsto |x|$, so that
$\Psi_t(\omega) = a(t)\omega^2 + \int_{(0,\infty)}(1 - \cos\omega v)\,\varpi_t(dv)$. Let $q > 1$.
Then (A8$|q^{\ZZ}$) holds with $S_qt = qt$ if and only if
\[
  \Psi_{qt}(\omega) = \Psi_t(q\omega) \ \text{ for all } t \ge 0,\ \omega \in \RR
  \qquad\Longleftrightarrow\qquad
  a(qt) = q^2a(t), \quad \varpi_{qt} = T_q\varpi_t \ \text{ for all } t > 0 .
\]
Every kernel is then semi-selfdecomposable with span $q$: $T_q\mu_{0,t} = \mu_{0,qt} =
\mu_{0,t} * \mu_{t,qt}$, and $T_q\varpi_t \ge \varpi_t$, $q^2a(t) \ge a(t)$. The family is
determined by its kernels $\mu_{0,t}$ at the scales $t$ of one \emph{fundamental annulus}
$[1, q)$.
\end{lemma}

\begin{proof}
By Lemma~\ref{lem:line-axioms-kernels} the kernels are a symmetric family in dimension $1$, and
(A8) at $q$ with $S_qt = qt$ is $q$-covariance of the kernels with that bijection. Conversely
$q$-covariance with $S_q t = qt$ gives (A8) at every $q^n$, $n \in \ZZ$, by the last sentence of
that lemma. So (A8$|q^{\ZZ}$) with $S_qt = qt$ is $q$-covariance of the kernels with
$S_qt = qt$.

For $0 \le s \le t$, (I1) gives $\widehat\mu_{0,t} = \widehat\mu_{s,t}\,\widehat\mu_{0,s}$, and
$\widehat\mu_{0,s} = e^{-\Psi_s}$ has no zero, so
$\widehat\mu_{s,t} = e^{-(\Psi_t - \Psi_s)}$. Covariance at the pairs $(0, t)$ reads
$\widehat\mu_{0,qt}(\omega) = \widehat\mu_{0,t}(q\omega)$, that is
$\Psi_{qt}(\omega) = \Psi_t(q\omega)$. Conversely, under this condition
$\widehat\mu_{qs,qt}(\omega) = e^{-(\Psi_t(q\omega) - \Psi_s(q\omega))} = \widehat\mu_{s,t}(q\omega)$,
the transform of $T_q\mu_{s,t}$, so $T_q\mu_{s,t} = \mu_{qs,qt}$ for all $s \le t$. For the
second equivalence, by Lemma~\ref{lem:line-exponent-data}(2) the function $\Psi_t(q\,\cdot)$ has
the folded data $(q^2a(t), T_q\varpi_t)$, and by (1) of that lemma a folded pair is determined by
its exponent. So $\Psi_{qt} = \Psi_t(q\,\cdot)$ if and only if $a(qt) = q^2a(t)$ and
$\varpi_{qt} = T_q\varpi_t$. At $t = 0$ both exponents vanish, since $\mu_{0,0} = \delta_0$.

Under covariance $T_q\mu_{0,t} = \mu_{0,qt}$, and (I1) gives
$\mu_{0,qt} = \mu_{t,qt} * \mu_{0,t}$. Reading this at $q^{-1}\omega$ gives
$\widehat\mu_{0,t}(\omega) = \widehat\mu_{0,t}(q^{-1}\omega)\,\widehat\varrho_q(\omega)$ with the
residual $\varrho_q := T_{q^{-1}}\mu_{t,qt}$. The step $\mu_{t,qt}$ is infinitely divisible by
Lemma~\ref{lem:family-structure} at $d = 1$, read through the embedding $x \mapsto xe_1$ of $\RR$
in $\RR^1$, and so is its dilate; so the kernel is semi-selfdecomposable with span $q$. The data
are nondecreasing: $\Psi_{qt} - \Psi_t$ is the exponent of the symmetric infinitely divisible law
$\mu_{t,qt}$, so it lies in $\mathrm{ND}_s$ (Lemma~\ref{lem:line-exponent-data}(1)), and
Lemma~\ref{lem:line-exponent-data}(3) gives $\varpi_t \le \varpi_{qt} = T_q\varpi_t$. Also
$q^2a(t) \ge a(t)$, since $a(t) \ge 0$ and $q > 1$.

Finally, two such families whose kernels agree on $[1, q)$ have, after the embedding of the line
in $\RR^1$, symmetric families in dimension $1$ that are $q$-covariant with $S_qt = qt$ and agree
there; by Lemma~\ref{lem:discrete-covariance-levy} at $d = 1$ they agree at every pair, and so do
the families on the line, since the embedding is injective on laws. The operators are the
convolutions with their kernels, so they agree too.

Every step is machine-checked, and the cited facts enter only through
Lemmas~\ref{lem:family-structure}, \ref{lem:line-exponent-data} and
\ref{lem:discrete-covariance-levy}.
\end{proof}

For processes, clause (1) of the next theorem is Sato's description of an additive process in law
by its triplets, with its converse [\sscite{sato1999levy}, Thms.~9.7--9.8, pp.~51--52], and
clause (3) is the analogue of his theorem that the marginals of the self-similar additive
processes are exactly the self-decomposable laws [\sscite{sato1999levy}, Thm.~16.1, pp.~99--100]. The theorem reads them for
the operator axioms on the line.

% shared with blueprint thm:lattice-line
\begin{theorem}[the three strata on the line]
\label{thm:lattice-line}

Let $(\Phi_{s,t})$ satisfy (A1)--(A7) and (ND), with kernels $\mu_{s,t}$, exponents $\Psi_t$ and
folded data $(a(t), \varpi_t)$ as in Lemma~\ref{lem:discrete-covariance-line}.
\begin{enumerate}
\item \emph{(Every family.)} The maps $t \mapsto a(t)$ and $t \mapsto \varpi_t$ are
nondecreasing, every step $\mu_{s,t}$ is a symmetric infinitely divisible law with the folded
data $(a(t) - a(s), \varpi_t - \varpi_s)$, and $t \mapsto \Psi_t(\omega)$ is continuous for each
$\omega$. Conversely, for $t \ge 0$ let numbers $a(t) \ge 0$ and measures $\varpi_t$ on
$(0,\infty)$ with $\int\min(1, v^2)\,\varpi_t(dv) < \infty$ be given, with $a(0) = 0$ and
$\varpi_0 = 0$. Suppose that whenever $s < t$, $a(s) \le a(t)$, $\varpi_s \le \varpi_t$ setwise and
$(a(s), \varpi_s) \ne (a(t), \varpi_t)$, and that
$t \mapsto a(t)\omega^2 + \int(1 - \cos\omega v)\,\varpi_t(dv)$ is continuous for each $\omega$.
Then these are the folded data of the kernels of a family with (A1)--(A7) and (ND).
\item \emph{(Discrete covariance.)} (A8$|q^{\ZZ}$) holds with $S_qt = qt$ if and only if
$a(qt) = q^2a(t)$ and $\varpi_{qt} = T_q\varpi_t$ for all $t > 0$. Every kernel $\mu_{0,t}$ is
then semi-selfdecomposable with span $q$, and the family is a \emph{pyramid}: its kernel at the
scale $q^nt$ is the $T_{q^n}$-image of its kernel at $t$.
\item \emph{(Full covariance.)} (A8) holds if and only if there is an increasing bijection $\chi$
of $[0,\infty)$ such that the relabelled family $\mu'_{s,t} := \mu_{\chi^{-1}(s),\chi^{-1}(t)}$ has
the folded data $a'(t) = \alpha t^2$ and $\varpi'_t(dv) = k(v/t)\,v^{-1}dv$ for all $t > 0$, with
$\alpha \ge 0$, $k \ge 0$ nonincreasing on $(0,\infty)$ and
$\int_0^\infty\min(1, x^2)\,k(x)\,x^{-1}dx < \infty$. Then $\Psi'_t(\omega) = F(t\omega)$ with
$F(\omega) := \alpha\omega^2 + \int_0^\infty(1 - \cos\omega x)\,k(x)\,x^{-1}dx$, the admissible
exponent (\ref{eq:admissible-exponent}) with data $(\alpha, k)$ of Theorem~\ref{thm:cited-A4}, and the relabelled family satisfies
(A8) with the bijections $S'_\lambda t = \lambda t$, $\lambda > 0$.
\end{enumerate}
\end{theorem}

\begin{proof}
(1) By Lemma~\ref{lem:line-axioms-kernels}(1) the kernels are a symmetric family in dimension
$1$. Every step $\mu_{s,t}$ is symmetric, and infinitely divisible by
Lemma~\ref{lem:family-structure} at $d = 1$, read through the embedding $x \mapsto xe_1$ of $\RR$
in $\RR^1$. By (I1) its transform is $e^{-(\Psi_t - \Psi_s)}$, so $\Psi_t - \Psi_s$ lies in
$\mathrm{ND}_s$ (Lemma~\ref{lem:line-exponent-data}(1)), and Lemma~\ref{lem:line-exponent-data}(3)
gives $a(s) \le a(t)$, $\varpi_s \le \varpi_t$, and the folded data
$(a(t) - a(s), \varpi_t - \varpi_s)$ of $\Psi_t - \Psi_s$, which are those of $\mu_{s,t}$.
Continuity of $\Psi_t(\omega) = -\log\widehat\mu_{0,t}(\omega)$ in $t$ is (I2), with
$\widehat\mu_{0,t} > 0$.

For the converse put $\Psi_t(\omega) := a(t)\omega^2 + \int(1 - \cos\omega v)\,\varpi_t(dv)$. For
$0 \le s \le t$ the hypotheses give $a(s) \le a(t)$ and $\varpi_s \le \varpi_t$, so
$\Psi_t - \Psi_s$ lies in $\mathrm{ND}_s$ (Lemma~\ref{lem:line-exponent-data}(3)), and by
Lemma~\ref{lem:line-exponent-data}(1) it is $-\log\widehat\mu_{s,t}$ for a symmetric law
$\mu_{s,t}$. These laws satisfy (I1): the transform of $\mu_{t,t}$ is $1$, that of $\delta_0$,
and $\mu_{s,t} * \mu_{r,s}$ and $\mu_{r,t}$ both have the transform $e^{-(\Psi_t - \Psi_r)}$; a law
is determined by its transform. They satisfy (I2), since
$(s,t) \mapsto e^{-(\Psi_t(\omega) - \Psi_s(\omega))}$ is continuous on $0 \le s \le t$ by the
continuity of $t \mapsto \Psi_t(\omega)$. They satisfy (ND): if $\mu_{s,t} = \delta_0$ for some
$s < t$, then $\Psi_s = \Psi_t$, and by the uniqueness of the folded pair, Theorem~\ref{thm:cited-A9},
$(a(s), \varpi_s) = (a(t), \varpi_t)$, against the hypothesis. So they form a symmetric family in
dimension $1$, and Lemma~\ref{lem:line-axioms-kernels}(2) turns it into a family of operators with
(A1)--(A7) and (ND) whose kernels are the $\mu_{s,t}$. Since $a(0) = 0$ and $\varpi_0 = 0$,
$\Psi_0 = 0$, so $\widehat\mu_{0,t} = e^{-\Psi_t}$. Thus $(a(t), \varpi_t)$ is a folded pair of
the exponent of $\mu_{0,t}$, and by the uniqueness of the data
(Lemma~\ref{lem:line-exponent-data}(1)) it is the folded data of $\mu_{0,t}$.

(2) The equivalence and the semi-selfdecomposability are
Lemma~\ref{lem:discrete-covariance-line}, and so is $T_{q^n}\mu_{0,t} = \mu_{0,q^nt}$, by
iterating $T_q\mu_{0,t} = \mu_{0,qt}$ in both directions.

(3) Suppose (A8). Theorem~\ref{thm:cited-A4}, applied to the family of operators, whose kernels
are the $\mu_{s,t}$ since an operator determines its kernel, gives an increasing bijection $\chi$
of $[0,\infty)$ and an exponent $F \not\equiv 0$,
$F(\omega) = \alpha\omega^2 + \int_0^\infty(1 - \cos\omega x)\,k(x)\,x^{-1}dx$ with $\alpha \ge 0$,
$k \ge 0$ nonincreasing and $\int\min(1,x^2)\,k(x)\,x^{-1}dx < \infty$, such that
$\widehat\mu_{s,t}(\omega) = \exp(-[F(\chi(t)\omega) - F(\chi(s)\omega)])$. At $s = 0$, where
$\chi(0) = 0$, $\chi$ being an increasing bijection of $[0,\infty)$, and $F(0) = 0$, this is
$\Psi_t(\omega) = F(\chi(t)\omega)$ for $t \ge 0$. So the relabelled family
$\mu'_{s,t} := \mu_{\chi^{-1}(s),\chi^{-1}(t)}$ has the exponents $\Psi'_t(\omega) = F(t\omega)$.
For $t > 0$ the substitution $v = tx$, with $dx/x = dv/v$, gives
\[
  F(t\omega) = \alpha\,t^2\omega^2 + \int_0^\infty(1 - \cos\omega v)\,k(v/t)\,v^{-1}dv ,
\]
and $k(\cdot/t)$ is again nonnegative and nonincreasing. The substitution gives
\[
  \int_0^\infty\min(1, v^2)\,k(v/t)\,\frac{dv}{v} = \int_0^\infty\min(1, t^2x^2)\,k(x)\,\frac{dx}{x}
  \le \max(1, t^2)\int_0^\infty\min(1, x^2)\,k(x)\,\frac{dx}{x} < \infty ,
\]
so this is the exponent of a folded pair. By the
uniqueness of the folded pair, Theorem~\ref{thm:cited-A9}, the data of $\mu'_{0,t}$ are $a'(t) = \alpha t^2$ and
$\varpi'_t(dv) = k(v/t)\,v^{-1}dv$.

Conversely, suppose that after a relabelling by an increasing bijection $\chi$ the data are of
this form. The relabelled kernels are again a symmetric family in dimension $1$: $\chi^{-1}$ is an
increasing bijection of $[0,\infty)$, hence an increasing homeomorphism fixing $0$, so it keeps the
cascade (I1), the continuity (I2), the inequality $s < t$ of (ND) and the laws themselves. By
Lemma~\ref{lem:family-structure} their transforms $\widehat\mu'_{0,s}$ are positive, and (I1) gives
$\widehat\mu'_{s,t} = e^{-(\Psi'_t - \Psi'_s)}$. Then $\Psi'_t(\omega) = F(t\omega)$ for $t > 0$,
by the substitution above read
backwards, and for $t = 0$, where both sides vanish since $\chi^{-1}(0) = 0$. So
$\Psi'_t(\lambda\omega) = \Psi'_{\lambda t}(\omega)$ for every $\lambda > 0$, and
\[
  \widehat{T_\lambda\mu'_{s,t}}(\omega) = \widehat\mu'_{s,t}(\lambda\omega)
  = e^{-(\Psi'_t(\lambda\omega) - \Psi'_s(\lambda\omega))}
  = e^{-(\Psi'_{\lambda t}(\omega) - \Psi'_{\lambda s}(\omega))}
  = \widehat\mu'_{\lambda s,\lambda t}(\omega) ,
\]
using $\widehat\mu'_{s,t} = e^{-(\Psi'_t - \Psi'_s)}$. So the relabelled kernels are
$\lambda$-covariant with $S'_\lambda t = \lambda t$ for every $\lambda > 0$, and by
Lemma~\ref{lem:line-axioms-kernels}(3), applied to the relabelled operators
$\Phi_{\chi^{-1}(s),\chi^{-1}(t)}$, whose kernels they are, these operators satisfy (A8) with the
bijections $S'_\lambda$. The original operators then satisfy (A8) with the increasing bijections
$S_\lambda := \chi^{-1}S'_\lambda\chi$ of $[0,\infty)$.

Every step is machine-checked, and the cited facts are \ledger{A4} and \ledger{A9} and those of
the lemmas named.
\end{proof}

Clause (1) is the first stratum, clause (2) the second and clause (3) the third. The two examples
below show them apart (Theorem~\ref{cor:lattice-line-strict}).

% shared with blueprint prop:lindeberg-stratum-one
\begin{proposition}[Lindeberg's discrete scale space lies strictly in the first stratum]
\label{prop:lindeberg-stratum-one}

Fix $h > 0$. For $0 \le s \le t$ let $\mu_{s,t}$ be the law on $h\ZZ$ with
$\mu_{s,t}(\{hn\}) = e^{-(t-s)}I_n(t-s)$, $n \in \ZZ$, where
$I_n(u) = \sum_{j \ge 0}(u/2)^{2j+|n|}/\bigl(j!\,(j+|n|)!\bigr)$ is the modified Bessel function of
the first kind. These are the kernels of the scale space for discrete signals of
[\sscite{lindeberg1990scalespace}], the semigroup generated on $h\ZZ$ by half the second difference.
\begin{enumerate}
\item $\widehat\mu_{s,t}(\omega) = e^{-(t-s)(1 - \cos h\omega)}$. So $\mu_{0,t}$ is compound
Poisson, with $\Psi_t(\omega) = t(1 - \cos h\omega)$, $a(t) = 0$ and $\varpi_t = t\,\delta_h$.
\item The operators $\Phi_{s,t}f := \mu_{s,t} * f$ satisfy (A1)--(A7) and (ND).
\item For no $q > 1$ and no $t > 0$ is $\mu_{0,t}$ semi-selfdecomposable with span $q$. So for no
$q > 1$ does the family satisfy (A8$|q^{\ZZ}$): it lies in the first stratum of
Theorem~\ref{thm:lattice-line} and not in the second.
\end{enumerate}
\end{proposition}

\begin{proof}
(1) Let $u \ge 0$ and $\theta \in \RR$. The two exponential series give
\[
  e^{u\cos\theta} = e^{(u/2)e^{i\theta}}\,e^{(u/2)e^{-i\theta}}
  = \sum_{a, b \ge 0}\frac{(u/2)^{a+b}}{a!\,b!}\,e^{i(a-b)\theta} ,
\]
the double series converging absolutely. The pairs $(a, b)$ with $a - b = n$ are
$(j + n, j)$, $j \ge 0$, when $n \ge 0$, and $(j, j + |n|)$ when $n < 0$; in both cases
$a + b = 2j + |n|$ and $a!\,b! = j!\,(j + |n|)!$. Grouping the terms by $n$ gives
\[
  \sum_{n \in \ZZ}I_n(u)\,e^{in\theta} = e^{u\cos\theta} .
\]
At $\theta = 0$ the masses $e^{-u}I_n(u) \ge 0$ sum to $1$, so each $\mu_{s,t}$ is a probability
law. At $\theta = h\omega$, with $u = t - s$,
$\widehat\mu_{s,t}(\omega) = \sum_n e^{-u}I_n(u)\,e^{ihn\omega} = e^{-u(1 - \cos h\omega)}$. So
$\Psi_t(\omega) = t(1 - \cos h\omega) = \int_{(0,\infty)}(1 - \cos\omega v)\,t\,\delta_h(dv)$.
The transform is real and even, so $\mu_{0,t}$ is symmetric, and it is the $n$-th power of the
transform of $\mu_{0,t/n}$ for every $n$, so $\mu_{0,t}$ is infinitely divisible. Its exponent has
the folded pair $(0, t\,\delta_h)$, and the folded data of a symmetric infinitely divisible law
are unique (Lemma~\ref{lem:line-exponent-data}(1)). So $\mu_{0,t}$ has the folded data
$a(t) = 0$ and $\varpi_t = t\,\delta_h$ (Definition~\ref{def:self-decomposable}), a
finite measure, and is compound Poisson. Once (2) has made the $\mu_{s,t}$ the kernels of a family,
these are the data $(a(t), \varpi_t)$ of Theorem~\ref{thm:lattice-line}.

(2) From the transforms, $\mu_{t,t} = \delta_0$ and $\mu_{s,t} * \mu_{r,s} = \mu_{r,t}$, which is
(I1). The transform $e^{-(t-s)(1 - \cos h\omega)}$ is continuous in $(s,t)$, which is (I2). It is
even in $\omega$, so every $\mu_{s,t}$ is symmetric, and for $s < t$ it is
$e^{-2(t-s)} \ne 1$ at $\omega = \pi/h$, so $\mu_{s,t} \ne \delta_0$. So the laws form a symmetric
family in dimension $1$, and Lemma~\ref{lem:line-axioms-kernels}(2) gives (A1)--(A7) and (ND).

(3) Let $q > 1$ and $t > 0$, and read $\mu_{0,t}$ on $\RR^1$ through the embedding
$x \mapsto xe_1$. Its image $m_t$ is infinitely divisible, by (2) and
Lemma~\ref{lem:family-structure} at $d = 1$, and has the transform
$e^{-t(1 - \cos h|\omega|)} = \exp\int(e^{i\langle\omega, x\rangle} - 1)\,t\sigma(dx)$ with
$\sigma := \tfrac12(\delta_{he_1} + \delta_{-he_1})$; since $\sigma$ is symmetric,
$(0, t\sigma, 0)$ is a generating triplet of $m_t$. The measure
$T_q(t\sigma) = \tfrac{t}{2}(\delta_{qhe_1} + \delta_{-qhe_1})$ gives the sphere $|x| = h$ the mass
$0$, and $t\sigma$ gives it $t > 0$, so $T_q(t\sigma) \ge t\sigma$ fails. By the criterion
Theorem~\ref{thm:cited-A44} (where $D_q = T_q$), $m_t$ is
not semi-selfdecomposable with span $q$. Neither is $\mu_{0,t}$: the image in $\RR^1$ of an
infinitely divisible residual of $\mu_{0,t}$ at $q$ would be one of $m_t$.

Suppose now that the family satisfied (A8$|q^{\ZZ}$) for some $q > 1$, with bijections $S$. By
Lemma~\ref{lem:line-axioms-kernels}(3) the kernels are $q$-covariant with the bijection $S_q$, and
so are their images in $\RR^1$, a symmetric family in dimension $1$. By
Lemma~\ref{lem:family-dilations}(3) an increasing bijection $\chi$ of $[0,\infty)$ relabels the
image family into one that is $q$-covariant with $S_qt = qt$. The relabelled family is again a
symmetric family in dimension $1$, $\chi^{-1}$ being an increasing homeomorphism of $[0,\infty)$
that fixes $0$ and the laws being unchanged. By
Lemma~\ref{lem:discrete-covariance-levy} its kernel at the scale $1$ is semi-selfdecomposable with
span $q$. Since $\chi(0) = 0$, that kernel is $m_u$ with $u := \chi^{-1}(1) > 0$, against the first
sentence. So no $q > 1$ has (A8$|q^{\ZZ}$), and by Theorem~\ref{thm:lattice-line}(1) the family lies
in the first stratum only.

Every step is machine-checked, and the cited facts are \ledger{A44} and those of the lemmas named.
\end{proof}

The grid spacing $h$ is a preferred length, which covariance forbids. The end of
Section~\ref{sec:line-theory} of the paper compares Lindeberg's generator with that of a fully
covariant family.

The next family is that of Proposition~\ref{prop:stratum-two-example} of the paper at $d = 1$. The
next proposition writes the position of a scale within its octave as $\zeta$, where
Section~\ref{ch:pyramid} of the paper writes $\theta$.

% shared with blueprint prop:discrete-gaussian-pyramid
\begin{proposition}[the geometric-jump pyramid lies strictly in the second stratum]
\label{prop:discrete-gaussian-pyramid}

Fix $q > 1$, $h > 0$ and weights $w_i \ge 0$, $i \in \ZZ$, nonincreasing in $i$ and not all zero,
with $\sum_{i < 0}w_iq^{2i} < \infty$ and $\sum_{i \ge 0}w_i < \infty$. For $t = q^{n+\zeta}$
with $n \in \ZZ$ and $\zeta \in [0,1)$ put $w_i(\zeta) := (1-\zeta)w_i + \zeta w_{i-1}$,
\[
  \varpi_t := \sum_{i \in \ZZ}w_i(\zeta)\,\delta_{hq^{i+n}}, \qquad
  \Psi_t(\omega) := \sum_{i \in \ZZ}w_i(\zeta)\bigl(1 - \cos(hq^{i+n}\omega)\bigr),
\]
and $\varpi_0 := 0$, $\Psi_0 := 0$.
\begin{enumerate}
\item The laws $\mu_{s,t}$ with $\widehat\mu_{s,t} = e^{-(\Psi_t - \Psi_s)}$ are the kernels of a
$q^{\ZZ}$-covariant family in the canonical discrete gauge, with the folded data $a(t) = 0$ and
$\varpi_t$, and $T_q\varpi_t \ge \varpi_t$ for every $t$.
\item No increasing relabelling of the scales makes the family fully covariant: it lies in the
second stratum of Theorem~\ref{thm:lattice-line} and not in the third.
\item No kernel $\mu_{0,t}$, $t > 0$, is carried by a lattice $\xi\ZZ$, $\xi > 0$.
\end{enumerate}
\end{proposition}

\begin{proof}
(1) \emph{The pairs.} For $t > 0$ write $t = q^{n+\zeta}$ with $n \in \ZZ$ and $\zeta \in [0,1)$.
Since $w_{i-1} \ge w_i \ge 0$, the weights satisfy $w_i \le w_i(\zeta) \le w_{i-1}$. With
$1 - \cos x \le 2\min(x^2, 1)$,
\[
  w_i(\zeta)\bigl(1 - \cos(hq^{i+n}\omega)\bigr) \le 2\,w_{i-1}\min\bigl((hq^n\omega)^2q^{2i}, 1\bigr),
\]
and for any $c \in \RR$
\[
  \sum_{i \in \ZZ}w_{i-1}\min(c^2q^{2i}, 1) \le c^2q^2\sum_{j < 0}w_jq^{2j} + \sum_{j \ge 0}w_j
  < \infty ,
\]
splitting at $j = i - 1 = 0$. So the series $\Psi_t(\omega)$ converges, and, with $c = hq^n$,
$\int\min(1, v^2)\,\varpi_t(dv) < \infty$. The measure $\varpi_t$ is carried by the points
$hq^j > 0$, so $(0, \varpi_t)$ is a folded pair with the exponent $\Psi_t$; at $t = 0$ it is
$(0, 0)$.

\emph{Monotone and distinct.} The weight of $\varpi_t$ on the point $hq^j$ is
$W_t(j) := w_{j-n}(\zeta)$, the points $hq^j$ being distinct. Let $0 < s \le t$, $s = q^{m+\zeta'}$.
If $m = n$ then $\zeta' \le \zeta$ and $W_s(j) \le W_t(j)$, since $w_i(\zeta)$ is nondecreasing in
$\zeta$. If $m < n$ then $W_s(j) \le w_{j-m-1} \le w_{j-n} \le W_t(j)$. So
$\varpi_s \le \varpi_t$ setwise, and trivially so for $s = 0$. Some step of the weights is strict,
$w_{i_0} < w_{i_0-1}$: otherwise $w$ is constant, and $\sum_{i \ge 0}w_i < \infty$ makes it zero.
For $0 < s < t$ the atom $j = i_0 + m$ separates the two measures. If $m = n$ then $\zeta' < \zeta$
and $W_s(j) = w_{i_0} + \zeta'(w_{i_0-1} - w_{i_0}) < W_t(j)$. If $m < n$ then
$W_s(j) = (1-\zeta')w_{i_0} + \zeta'w_{i_0-1} < w_{i_0-1} \le w_{i_0+m-n} \le W_t(j)$, since
$\zeta' < 1$. For $s = 0 < t$, some $w_i > 0$, and $W_t(i + n) \ge w_i > 0$ while $\varpi_0 = 0$.
So $\varpi_s \ne \varpi_t$ whenever $s < t$.

\emph{Continuity.} Put $g_n(\omega) := \sum_iw_i\bigl(1 - \cos(hq^{i+n}\omega)\bigr)$. On the closed
octave $[q^n, q^{n+1}]$, with $\zeta := \log_qt - n$,
\[
  \Psi_t(\omega) = (1 - \zeta)\,g_n(\omega) + \zeta\,g_{n+1}(\omega) :
\]
for $t < q^{n+1}$ split the series along $w_i(\zeta) = (1-\zeta)w_i + \zeta w_{i-1}$ and shift the index
of the second part by one; at $t = q^{n+1}$ both sides are $g_{n+1}(\omega)$. So
$t \mapsto \Psi_t(\omega)$ is continuous on each closed octave, hence on $(0,\infty)$. Further
$g_n(\omega) = G(q^n)$ with $G(c) := \sum_iw_i\bigl(1 - \cos(hq^ic\omega)\bigr)$, a series of
continuous functions dominated on $[-1,1]$ by the convergent series
$2\sum_iw_{i-1}\min((h\omega)^2q^{2i}, 1)$, so $G$ is continuous there, with $G(0) = 0$. Hence
$g_n(\omega) \to 0$ as $n \to -\infty$, and since $0 \le \Psi_t(\omega) \le \max(g_n(\omega),
g_{n+1}(\omega))$ on the octave of $t$, $\Psi_t(\omega) \to 0 = \Psi_0(\omega)$ as $t \downarrow 0$.

\emph{The family.} The pairs $(0, \varpi_t)$ thus satisfy the hypotheses of the converse in
Theorem~\ref{thm:lattice-line}(1), which gives a family with (A1)--(A7) and (ND) whose kernels
$\mu^0_{s,t}$ have the folded data $(0, \varpi_t)$. By (I1) their transforms are
$\widehat\mu^0_{s,t} = e^{-(\Psi_t - \Psi_s)}$, so laws with these transforms exist, and any such laws
are the $\mu^0_{s,t}$, a law being determined by its transform. They are therefore the kernels of
that family, with the folded data $(0, \varpi_t)$. For $t > 0$, $qt = q^{(n+1)+\zeta}$, so
$\Psi_{qt}(\omega) = \Psi_t(q\omega)$, and $\Psi_0 = 0$. By Lemma~\ref{lem:discrete-covariance-line}
the family satisfies (A8$|q^{\ZZ}$) with $S_qt = qt$, so it is $q^{\ZZ}$-covariant in the canonical
discrete gauge (Lemma~\ref{lem:line-axioms-kernels}), and $T_q\varpi_t \ge \varpi_t$. Directly:
$T_q\varpi_t$ gives the point $hq^{i+n}$ the weight $w_{i-1}(\zeta) \ge w_i(\zeta)$.

(2) Let operators $\Phi_{s,t}$ have the kernels $\mu_{s,t}$. An operator with a kernel is
convolution by it, so they are the operators of (1). Suppose an increasing bijection $\psi$ of
$[0,\infty)$ made the relabelled operators $\Phi_{\psi(s),\psi(t)}$ satisfy (A8), with bijections
$S_\lambda$, and let $\chi := \psi^{-1}$. Then the family itself satisfies (A8), with the increasing
bijections $\psi \circ S_\lambda \circ \chi$. By Theorem~\ref{thm:lattice-line}(3) there are an
increasing bijection $\chi'$ of $[0,\infty)$ and data $(\alpha, k)$ such that the relabelled family
$\mu_{\psi'(s),\psi'(t)}$, $\psi' := \chi'^{-1}$, has the folded data $\varpi_{\psi'(t)}(dv) =
k(v/t)\,v^{-1}dv$ for $t > 0$, a measure with a density, which gives every point the mass $0$. But
$u := \psi'(1) > 0$, and $\varpi_u$ gives some point $hq^j$ a positive mass, as shown in (1). So no
such $\psi$ exists.

(3) The family is $q^{\ZZ}$-covariant, and this is
Proposition~\ref{prop:no-lattice-kernel-line}(1).

Every step is machine-checked, and the cited facts enter only through
Theorem~\ref{thm:lattice-line}.
\end{proof}

The family repairs Lindeberg's discrete scale space for covariance. Lindeberg's kernels, his
discrete analogue of the Gaussian, jump by the one length $h$. Here each octave adds the jumps of
the next length $hq^i$. The lengths of the jumps form a geometric sequence, hence the name. The
family has no Gaussian part, and its \Levy{}
measures are atomic. In $\RR^d$, $d \ge 2$, no rotation-invariant infinitely divisible law has a
marginal whose folded \Levy{} measure has an atom (Proposition~\ref{prop:what-does-not-lift}), so
the examples of Section~\ref{sec:pyramid-strata} of the paper put their jumps on spheres.

% shared with blueprint cor:lattice-line-strict
\begin{theorem}[the three strata are distinct]
\label{cor:lattice-line-strict}

Fix $q > 1$. As conditions on families with (A1)--(A7) and (ND), full covariance implies
(A8$|q^{\ZZ}$), and both implications of the chain
\emph{fully covariant} $\Rightarrow$ \emph{$q^{\ZZ}$-covariant} $\Rightarrow$ \emph{(A1)--(A7) and
(ND)} are strict: Lindeberg's discrete scale space
(Proposition~\ref{prop:lindeberg-stratum-one}) is in the first stratum of
Theorem~\ref{thm:lattice-line} and not in the second, and the geometric-jump pyramid
(Proposition~\ref{prop:discrete-gaussian-pyramid}) is in the second and not in the third.
\end{theorem}

\begin{proof}
(A8) requires the covariance at every $\lambda > 0$, so in particular at every $\lambda \in q^{\ZZ}$,
which is (A8$|q^{\ZZ}$). Lindeberg's discrete scale space, at any $h > 0$, satisfies (A1)--(A7)
and (ND) and not (A8$|q^{\ZZ}$) (Proposition~\ref{prop:lindeberg-stratum-one}(2)--(3)). The
geometric-jump pyramid satisfies (A8$|q^{\ZZ}$), and no relabelling of its scales makes it fully
covariant (Proposition~\ref{prop:discrete-gaussian-pyramid}(1)--(2)). It is not fully covariant in
its own labelling either, the identity being a relabelling. For the given $q$ one instance is
$h = 1$ with the weights $w_i = 1$ for $i < 0$ and $w_i = 0$ for $i \ge 0$: they are nonnegative,
nonincreasing and not all zero, $\sum_{i<0}q^{2i} < \infty$, and $\sum_{i \ge 0}w_i = 0$.
\end{proof}

% shared with blueprint prop:no-lattice-kernel-line
\begin{proposition}[no kernel from scale zero is a lattice law, and the pyramid principle]
\label{prop:no-lattice-kernel-line}

Let $q > 1$ and let $(\Phi_{s,t})$ be $q^{\ZZ}$-covariant, with kernels $\mu_{s,t}$.
\begin{enumerate}
\item No kernel $\mu_{0,t}$ with $t > 0$ is carried by a lattice $\xi\ZZ$, $\xi > 0$.
\item Suppose $S_qt = qt$, and that $\mu_{0,t}$ has a continuous density $\phi_t$ for some
$t > 0$. Then $\mu_{0,q^nt}$ has the continuous density $\phi_{q^nt}(x) = q^{-n}\phi_t(q^{-n}x)$
for every $n \in \ZZ$, so $q^n\phi_{q^nt}(q^nhm) = \phi_t(hm)$ for every $h > 0$ and $m \in \ZZ$:
sampled at the scale $q^nt$ on the grid $q^nh\ZZ$, the kernel is the same array at every level,
up to the factor $q^n$.
\end{enumerate}
\end{proposition}

\begin{proof}
(1) By Lemma~\ref{lem:line-axioms-kernels} the kernels form a symmetric family in dimension $1$
that is $q$-covariant. Its image in $\RR^1$ under $x \mapsto xe_1$ is a symmetric $q$-covariant
family in the sense of Definition~\ref{def:family-dimension-d}, and the lattice $\xi\ZZ$ goes to
the subgroup generated by $\xi e_1$, which is discrete: its nonzero points have norm at least
$\xi$. Proposition~\ref{prop:no-lattice-kernel} at $d = 1$ says that no kernel of the image, at
$t > 0$, is carried by it, and the image of $\mu_{0,t}$ is carried by it exactly when $\mu_{0,t}$ is
carried by $\xi\ZZ$.

(2) In the canonical gauge, $T_q\mu_{0,t} = \mu_{0,qt}$ for every $t \ge 0$
(Lemma~\ref{lem:discrete-covariance-line}), and iterating this and its inverse,
$\mu_{0,q^nt} = T_{q^n}\mu_{0,t}$ for every $n \in \ZZ$. If $\mu$ has the density $\phi$, then
$T_c\mu$ has the density $c^{-1}\phi(x/c)$, by the substitution $x = cy$. This density is
continuous. Two continuous densities of one law agree almost everywhere, hence everywhere, since a
nonempty open set has positive Lebesgue measure. So it is the continuous density of
$\mu_{0,q^nt}$, and at $x = q^nhm$ it is $q^{-n}\phi_t(hm)$.
\end{proof}

A step $\mu_{s,t}$ with $s > 0$ can be a lattice law, and that is what a grid pyramid is
(Definition~\ref{def:grid-pyramid-line}, Proposition~\ref{prop:grid-pyramids-infinitely-divisible}(4)).
Clause (2) is the pyramid principle of Section~\ref{subsec:line-strata} of the paper. The kernels
from scale zero of a $q^{\ZZ}$-covariant family are carried by no lattice, by (1). When such a kernel
has a continuous density, covariance in the canonical gauge makes its samples on the grids
$q^nh\ZZ$ one array.
